%==============================================================================
%  WEEKLY AI BRIEF -- issue #04, window 9-15 September 2026
%  Engine: tectonic (also compiles with pdflatex/latexmk)
%==============================================================================
\documentclass[11pt,a4paper]{article}

%------------------------------ packages -------------------------------------
\usepackage[T1]{fontenc}
\usepackage[utf8]{inputenc}
\usepackage{lmodern}
\usepackage[margin=2.2cm,top=2.0cm,bottom=2.3cm]{geometry}
\usepackage{microtype}
\usepackage{graphicx}
\usepackage{xcolor}
\usepackage{titlesec}
\usepackage{enumitem}
\usepackage{booktabs}
\usepackage{tabularx}
\usepackage{fancyhdr}
\usepackage{lastpage}
\usepackage{tcolorbox}
\usepackage{needspace}
\usepackage[hidelinks]{hyperref}

\tcbuselibrary{skins,breakable}

% Logos live in template/ and are shared by every issue -- never copy them into a
% dated folder. This path list resolves them from <YYYY/MM/DD>/ or from template/.
\graphicspath{{../../../template/}{../template/}{./}{template/}}

%------------------------------ palette --------------------------------------
\definecolor{brandink}{HTML}{103868}    % KAUST navy -- headings, rules
\definecolor{brandaccent}{HTML}{0B7480} % KAUST teal (darkened to 5.6:1 on white)
\definecolor{brandsoft}{HTML}{EAF3F5}   % box fills
\definecolor{brandgrey}{HTML}{5A6672}   % meta text

\hypersetup{colorlinks=true,linkcolor=brandaccent,urlcolor=brandaccent,
            citecolor=brandaccent,pdfborder={0 0 0}}

%------------------------------ metadata -------------------------------------
\newcommand{\briefnumber}{04}
\newcommand{\briefweek}{9--15 September 2026}
\newcommand{\briefdate}{15 September 2026}
\newcommand{\briefauthor}{Ali Habibullah}
\newcommand{\briefaffil}{KAUST}
\newcommand{\briefscope}{Papers, model releases, benchmarks and policy from the past seven days}

%------------------------------ section styling ------------------------------
\titleformat{\section}
  {\normalfont\sffamily\large\bfseries\color{brandink}}{\thesection.}{0.6em}{}
  [\vspace{-0.55em}{\color{brandaccent}\rule{\textwidth}{0.9pt}}]
\titlespacing*{\section}{0pt}{1.5ex plus .3ex}{0.9ex}

\titleformat{\subsection}[runin]
  {\normalfont\sffamily\bfseries\color{brandink}}{}{0pt}{}[\;---]
\titlespacing*{\subsection}{0pt}{1.25ex plus .2ex}{0.5em}

\setlist[itemize]{leftmargin=1.15em,itemsep=0.25em,topsep=0.35em,parsep=0pt}
\setlength{\parindent}{0pt}
\setlength{\parskip}{0.45em}

%------------------------------ header / footer ------------------------------
\pagestyle{fancy}
\fancyhf{}
\renewcommand{\headrulewidth}{0pt}
\renewcommand{\footrulewidth}{0.4pt}
\fancyfoot[L]{\footnotesize\color{brandgrey}Weekly AI Brief \#\briefnumber\ \textbullet\ \briefweek}
\fancyfoot[R]{\footnotesize\color{brandgrey}Page \thepage\ of \pageref{LastPage}}

%------------------------------ content macros -------------------------------
\newcommand{\tldr}[1]{%
  \begin{tcolorbox}[breakable,enhanced,colback=brandsoft,colframe=brandaccent,
      boxrule=0.8pt,left=10pt,right=10pt,top=8pt,bottom=8pt,arc=2pt,
      title={\sffamily\bfseries The week in five lines},fonttitle=\color{white}]
    \begin{itemize}[leftmargin=1.1em,itemsep=0.3em,topsep=0pt]#1\end{itemize}
  \end{tcolorbox}}

\newcommand{\entry}[4]{%
  \par\vspace{0.7em}%
  \Needspace*{5\baselineskip}%
  {\sffamily\bfseries\color{brandink}#1}\par
  {\footnotesize\color{brandgrey}#2 \textbullet\ \href{#3}{link} \textbullet\ \srcref{#4}}\par
  \vspace{0.15em}}

\newcommand{\why}[1]{\par{\small\textbf{Why it matters.} #1}\par}

\newcommand{\srcref}[1]{\textcolor{brandaccent}{[S#1]}}

\newcommand{\stat}[2]{{\sffamily\bfseries\color{brandink}#1}~{\footnotesize\color{brandgrey}#2}}

%==============================================================================
\begin{document}

%------------------------------ title block ----------------------------------
\noindent
\begin{minipage}[c]{0.48\textwidth}
  \raggedright\includegraphics[height=1.4cm]{KAUST_Academy_Logo.png}
\end{minipage}\hfill
\begin{minipage}[c]{0.48\textwidth}
  \raggedleft\includegraphics[height=1.5cm]{KAUST_Logo.png}
\end{minipage}
\vspace{0.7em}

{\sffamily
 {\color{brandgrey}\footnotesize ISSUE \#\briefnumber\ \textbullet\ \briefdate}\par
 \vspace{0.15em}
 {\LARGE\bfseries\color{brandink}Weekly AI Brief}\par
 \vspace{0.25em}
 {\large\color{brandaccent}\briefweek}\par
 \vspace{0.35em}
 {\footnotesize\color{brandgrey}\briefscope\\
  Prepared by \briefauthor, \briefaffil}\par}
\vspace{0.3em}
{\color{brandink}\rule{\textwidth}{1.4pt}}
\vspace{0.6em}

%------------------------------ 1. executive summary -------------------------
\tldr{
  \item \textbf{OpenAI says an internal model produced a finite-time blowup proof for
        forced Navier--Stokes} --- 166 pages plus a Lean formalization, from roughly 10{,}000
        agents in 88 hours. The Clay Institute says the problem has ``apparently been
        settled'' and that its review is ``deliberately unhurried''; a priority dispute with
        two mathematicians is unresolved. \srcref{1}\,\srcref{5}
  \item \textbf{DeepSeek-V4.1-Flash ships under MIT}: 552B-parameter MoE with 8B/16B active,
        1M context, \$0.30 in / \$1.20 out per million at peak. Artificial Analysis scores it
        40 on its index, five points above the previous Flash. \srcref{7}\,\srcref{9}
  \item \textbf{OpenAI turns its harness into products}: GPT-Live-1 reaches the API at
        \$0.05 per minute, and the Agents API opens in public beta with no fee beyond
        tokens. \srcref{11}\,\srcref{12}
  \item \textbf{An open recipe for IMO gold}: Nemotron checkpoints score 30 of 42 at IMO
        2026 in natural language, with training data, code and a 200-problem benchmark
        released. \srcref{14}
  \item \textbf{Anthropic's threat report covers seven harm domains} over nine months; none
        of the cases used Fable- or Mythos-class models, except one illicit
        distillation. \srcref{17}
}

%------------------------------ 2. the story of the week ---------------------
\section{The Navier--Stokes claim, and the dispute around it}

Three documents landed within four days and none of them settles the question of who did
what first. The OpenAI announcement is dated 8 September --- one day before this window
--- and is carried here because issue \#03 closed before it appeared, because OpenAI
amended it on 10 September, and because the Clay Institute answered on 11 September.

\entry{OpenAI: ``Finite time blowup for Navier--Stokes''}{OpenAI \textbullet\ 8 September, updated 10 September}{https://openai.com/index/navier-stokes-solution/}{1}
OpenAI states that an internal model ``significantly more capable than GPT-6 Astra'' produced
a proof that a three-dimensional incompressible fluid, starting from rest under a smooth,
compactly supported force, develops unbounded velocity in finite time while its kinetic
energy stays bounded --- statements ``C'' and ``D'' of the official Clay formulation, i.e.\ the
\emph{forced} blowup route, not the unforced regularity question. The write-up is
\textbf{166 pages} \srcref{2} and a Lean formalization is on GitHub. By OpenAI's account the
effort began on 1 September after rumours that two Millennium problems had fallen; a group
``on the order of 10{,}000 concurrent agents'' reached the result on 5 September, 88 hours
after launch, Lean verification took a further 17 hours ``via GPT-6 Astra'', and the
Navier--Stokes run alone used 2.7 million agent messages and about 130 billion output
tokens. A separate group of ``nearly 100 agents'' produced an unforced Euler blowup in about
50 hours. OpenAI says it does not intend to claim the prize. The 10 September update adds
that an internal investigation found the Codex prompts of Tristan Buckmaster ``could not have
influenced the system in any way, including through training''.
\why{Whatever the review concludes, the resource figures are the first public accounting
of what a frontier-scale multi-agent mathematics run costs in messages and tokens.}

\entry{Buckmaster and Alp\"oge: Euler blowup with smooth forcing, and a statement}{Buckmaster, T. (NYU) and Alp\"oge, L. (Anthropic) \textbullet\ 8 September}{https://cims.nyu.edu/~tristanb/statement.pdf}{3}
The same day, Buckmaster and Alp\"oge released three Lean-verified results --- finite-time
blowup with smooth forcing for incompressible porous media, 2D Boussinesq and 3D
incompressible Euler; the Euler paper runs to \textbf{112 pages} \srcref{4}. Buckmaster
credits the underlying program to Diego C\'ordoba and Luis Mart\'inez-Zoroa, dates the Euler
and Boussinesq results to 15 August with Lean verification on 22 August, and says the pair
used ``Anthropic's Claude, OpenAI's Codex, especially with GPT-5.6 Sol and, more recently,
Astra''. His statement then recounts a 3 September email to OpenAI and two 6 September calls
with S\'ebastien Bubeck in which he was told OpenAI's model had a roughly 100-page forced
Navier--Stokes proof, and was offered a sequenced joint release or sole authorship of a
paper on OpenAI's result without Alp\"oge. He declined both. He also writes: ``I have not seen
OpenAI's proof \ldots I am not accusing anyone of anything. I am stating what I was told,
when, and what was proposed to me.''
\why{Two AI-assisted proof programs converged on the same forced-blowup route within
days of each other, and the record of who prompted what, and when, is now part of the
mathematics.}

\entry{Clay Mathematics Institute statement}{CMI \textbullet\ 11 September}{https://www.claymath.org/news/navier-stokes-announcement/}{5}
The Institute ``shares in the excitement \ldots as we contemplate the announcement that the
Navier--Stokes problem has apparently been settled'', names no author, and says the rules
``describe the process for evaluating what has been achieved and for assigning credit. The
process is deliberately unhurried, but we will provide updates.'' None of the sources read for
this issue records an independent mathematician endorsing the OpenAI proof; the Lean
artifact is the only verification on record.
\why{``Apparently settled'' from the prize's custodian is as far as anyone has gone; the
review, not the announcement, is the event to wait for.}

%------------------------------ 3. models & releases -------------------------
\section{Model and product releases}

\entry{DeepSeek-V4.1-Flash}{DeepSeek \textbullet\ open weights, MIT \textbullet\ 10 September}{https://huggingface.co/deepseek-ai/DeepSeek-V4.1-Flash}{6}
DeepSeek's new ``Causal Encoder--Decoder'' design: a 40-layer transformer
split into a 20-layer causal encoder and a 20-layer decoder, with \textbf{552B} backbone
parameters and just \textbf{8B active per token in prefill and 16B in decode}. Context is
1M tokens, pre-training used 45T multimodal tokens, image input is native, and Compressed
Sparse Attention 2 brings the global KV cache to 890 bytes per token, about a quarter of
V4-Flash \srcref{7}. API pricing at peak is \$0.30 per million input tokens (cache miss),
\$1.20 output and \$0.006 on cache hits, with off-peak rates at half; max output is 384K
\srcref{8}. Self-reported at maximum reasoning effort: Terminal-Bench 2.1 \textbf{90.6\%}
against 89.1\% for Claude Opus 5 and 88.8\% for GPT-5.6 Sol, DeepSWE v1.1 74.2\%
(74.0\% / 73.0\%), but GPQA Diamond 90.9\% (93.4\% / 94.1\%), HLE 36.8\% (56.3\% / 44.5\%)
and Terminal-Bench 4.0 31.2\% against Opus 5's 51.8\% \srcref{7}. DeepSeek is retiring
V4-Pro: from 14 September its requests route to V4.1-Flash \srcref{6}.
\why{A 16B-active open model that trades blows with closed frontier models on agentic
coding at \$1.20 per million output tokens resets the baseline for self-hosted agents ---
on the vendor's own numbers.}

\entry{GPT-Live-1 in the API}{OpenAI \textbullet\ 10 September}{https://openai.com/index/introducing-gpt-live-1-in-the-api/}{11}
The full-duplex voice model that has powered ChatGPT Voice since July is now an API product
at \textbf{\$0.05 per minute} for the voice layer, with the backend reasoning model billed
separately. It listens and speaks in one model, delegates reasoning and tool calls to a
text model (GPT-6 Astra or a third-party model), supports telephony, and exposes ASR
transcripts, keyword biasing and native turn detection. OpenAI reports a 30-point gain on
Full Duplex Bench over GPT-Realtime-2.1 and first place on Tau3 when paired with Astra at
medium effort --- both self-reported.
\why{Pricing the voice layer per minute and the intelligence per token separates the two
costs cleanly, which is how voice agents will now be budgeted.}

\entry{Agents API (public beta)}{OpenAI \textbullet\ 10 September}{https://openai.com/index/introducing-the-agents-api/}{12}
The harness behind Codex as a managed service: one call specifies task, model, tools and
environment; OpenAI runs session orchestration, context compaction, recovery and subagent
fan-out. Compute can be an OpenAI-hosted sandbox, your own VPC, or a partner (Blaxel,
Cloudflare, Daytona, DigitalOcean, E2B, Modal, Oracle, Runloop, Vercel). No fee beyond
tokens and tools; the harness is the open-source Codex harness. Customer figures on the
page (a 60\% cost-per-case cut) are testimonials, not benchmarks.
\why{The harness is now a product surface --- and the Anthropic report below shows the
same pattern is what attackers build too.}


%------------------------------ 4. research ----------------------------------
\section{Research worth reading}

\entry{NCP-ArchPreview: Next Concept Prediction}{Intern-NCP Team \textbullet\ arXiv:2609.10715}{https://arxiv.org/abs/2609.10715}{13}
An 8.9B model trained on 5.73T Dolma-3 tokens that predicts discrete multi-token
``concepts'' (product-quantized from hidden states) alongside the next token. The authors
report reaching OLMo-3-7B's final pre-training loss using \textbf{51.3\%} of its tokens, a
2.45-point average downstream gain and +5.99 on GSM8K. Submitted 9 September; the abstract
does not state whether weights are released.
\why{A halving of token budget at matched loss is the kind of claim worth reproducing
before anyone changes a pre-training recipe on the strength of it.}

\entry{An Open Recipe for IMO Gold: Training Nemotron for Olympiad Mathematics}{Moshkov, I. et al. \textbullet\ arXiv:2609.10712}{https://arxiv.org/abs/2609.10712}{14}
Starting from Nemotron 3 Ultra, SFT-then-RL specialist checkpoints score \textbf{30 of 42}
at IMO 2026 --- gold --- in natural language with no formal prover or tools. Released: two
checkpoints, training data, training and inference code, the submitted solutions, and
Nemotron-IMO-Bench, 200 new olympiad problems. Submitted 9 September.
\why{Every prior IMO-gold result came from a closed lab; this is the first with the full
recipe on the table.}

\entry{ZGCM-1: a fully open 7B model for math and agentic search}{He, J. et al. \textbullet\ arXiv:2609.13356 \textbullet\ CC BY 4.0}{https://arxiv.org/abs/2609.13356}{15}
A 7B model trained from scratch with a 256K context, interleaved gated sliding-window and
full attention, and an FP8 Muon optimizer, reporting about \textbf{4.2$\times$} better
time-to-loss at 16K pre-training. What sets it apart is the release: weights from
pre-, mid- and post-training, intermediate checkpoints, training code, per-stage data and
recipes, and the W\&B logs. Submitted 11 September.
\why{Complete training logs and per-stage data are what make an ``open'' model
auditable; almost no release at this scale includes them.}


%------------------------------ 5. benchmarks & evals ------------------------
\section{Benchmarks, evaluations and reproductions}

One independent measurement and one benchmark repair this week. Artificial Analysis
evaluated DeepSeek-V4.1-Flash on its own harness \srcref{9}, and the same index version
scores the previous Flash \srcref{10}. Everything else in the table is self-reported.

{\footnotesize
\begin{tabularx}{\textwidth}{@{}l l r X@{}}
\toprule
\textbf{Benchmark} & \textbf{Model} & \textbf{Score} & \textbf{Who ran it} \\
\midrule
AA Intelligence Index v4.3 & DeepSeek-V4.1-Flash (max) & 40 & Artificial Analysis; V4 Flash 0731 scores 35 on the same index \srcref{9}\,\srcref{10} \\
Terminal-Bench 2.1        & DeepSeek-V4.1-Flash & 90.6\% & DeepSeek; vs 89.1\% Opus 5, 88.8\% Sol \srcref{7} \\
Terminal-Bench 4.0        & DeepSeek-V4.1-Flash & 31.2\% & DeepSeek; vs 51.8\% Opus 5, 39.9\% Sol \srcref{7} \\
Full Duplex Bench         & GPT-Live-1 & +30 pts & OpenAI, over GPT-Realtime-2.1 \srcref{11} \\
IMO 2026 (natural language) & Nemotron IMO checkpoints & 30/42 & Authors; artifacts released \srcref{14} \\
\bottomrule
\end{tabularx}}
\vspace{0.4em}

\textbf{SWE-Bench Pro Verified} \srcref{16} (arXiv:2609.08149; submitted 8 September,
surfaced on 10 September) documents reward hacking on SWE-Bench Pro through leaked gold
solutions and hidden evaluation information, plus misleading problem statements and
badly scoped tests, and ships a repaired set with anti-leakage safeguards. The authors
report that ``some models perform substantially worse than previously reported''. The
abstract gives no per-model numbers and no artifact link, so treat it as a warning about the
existing leaderboard rather than a new one.

%------------------------------ 6. industry & policy -------------------------
\section{Industry, funding and policy}

\begin{itemize}
  \item \textbf{Anthropic's ``Detecting and countering misuse of AI: September 2026''}
        (10 September) covers December 2025 to August 2026 across seven domains: cyber,
        surveillance, influence, conventional weapons, biological misuse, scams and fraud,
        and illicit distillation. Named cases include a Russian-speaking operator linked to
        Midnight Blizzard with 300{,}000+ identity records stolen from 20+ organisations;
        a Chinese-speaking group with about fifty targets and ``a dozen possible zero day
        findings in a single month''; and a France-based agency running about 70
        fabricated news sites with 8{,}913 articles in about 20 languages. ``None of the misuse cases involved the use of Claude Fable
        or Mythos-class models, with the exception of one illicit distillation case.''
        \srcref{17}
  \item \textbf{Paul Christiano joins the OpenAI Foundation Board} (9 September) and its
        Safety and Security Committee under chair Zico Kolter, as a non-voting observer on
        the OpenAI Group PBC board. He remains Senior Technical Advisor at NIST's CAISI and
        will recuse himself from OpenAI-related matters and model evaluations there.
        \srcref{18}
  \item \textbf{Gulf compute:} Fortune reports (9 September), relaying Bloomberg, that
        Humain is seeking a \$2.5 billion fund for Saudi data-centre expansion; the piece
        puts Saudi capacity at 467 MW in Q1 2026 against a Humain target of 1.9 GW by 2030,
        and CEO Tareq Amin's stated plan to list in Riyadh and New York by 2029. A
        reported fundraising, not a closed one. \srcref{19}
  \item \textbf{Missed by issues \#02 and \#03:} NVIDIA announced on 3 September that it
        will acquire Hugging Face for \$12{,}930{,}300{,}000, pledging that the platform
        ``will remain an open platform for the entire AI ecosystem'' and that ``NVIDIA
        compute will not be required''. Outside this window; included because the brief has
        not yet recorded it. \srcref{20}
\end{itemize}

%------------------------------ 7. what to watch -----------------------------
\section{What to watch next week}
\begin{itemize}
  \item \textbf{Independent reading of the 166-page proof.} The Clay Institute has
        promised updates \srcref{5}; the first named analyst to endorse or fault the
        argument, rather than the Lean certificate, will decide how this is remembered.
  \item \textbf{DeepSeek-V4.1-Flash on neutral harnesses.} Artificial Analysis has one
        number so far \srcref{9}; third-party Terminal-Bench 4.0 and SWE-style runs will
        show whether the 16B-active claim survives outside DeepSeek's own evaluation.
  \item \textbf{Whether the SWE-Bench Pro repair changes any published score.} If the
        leakage the authors describe is real \srcref{16}, vendor coding claims made on the
        unrepaired set need re-running.
\end{itemize}

%------------------------------ 8. sources -----------------------------------
\Needspace*{10\baselineskip}
\section{Sources}
{\footnotesize
\begin{tabularx}{\textwidth}{@{}l X@{}}
\toprule
\textbf{\#} & \textbf{Reference} \\
\midrule
S1 & OpenAI. \emph{On the Navier--Stokes Millennium Prize Problem}. 8 September 2026, updated 10 September. \url{https://openai.com/index/navier-stokes-solution/} \\
S2 & OpenAI. \emph{Finite time blowup for Navier--Stokes} (166 pp.). \url{https://cdn.openai.com/pdf/32d9f210-8b73-45e0-91bc-82a30aef8a9a/navier-stokes.pdf} \\
S3 & Buckmaster, T. \emph{Statement}. 8 September 2026. \url{https://cims.nyu.edu/~tristanb/statement.pdf} \\
S4 & Alp\"oge, L. and Buckmaster, T. \emph{Blowup for the Euler equations with smooth forcing} (112 pp.). \url{https://cims.nyu.edu/~tristanb/euler.pdf} \\
S5 & Clay Mathematics Institute. \emph{Navier--Stokes announcement}. 11 September 2026. \url{https://www.claymath.org/news/navier-stokes-announcement/} \\
S6 & DeepSeek. \emph{Introducing DeepSeek-V4.1-Flash}. 10 September 2026. \url{https://www.deepseek.com/en/news/deepseek-v4-1-flash/} \\
S7 & DeepSeek. \emph{DeepSeek-V4.1-Flash model card}. Accessed 15 September 2026. \url{https://huggingface.co/deepseek-ai/DeepSeek-V4.1-Flash} \\
S8 & DeepSeek. \emph{Models \& Pricing}. Accessed 15 September 2026. \url{https://api-docs.deepseek.com/quick_start/pricing} \\
S9 & Artificial Analysis. \emph{DeepSeek V4.1 Flash (max)}. Accessed 15 September 2026. \url{https://artificialanalysis.ai/models/deepseek-v4-1-flash} \\
S10 & Artificial Analysis. \emph{DeepSeek V4 Flash 0731 (max)}. Accessed 15 September 2026. \url{https://artificialanalysis.ai/models/deepseek-v4-flash} \\
S11 & OpenAI. \emph{Build more natural voice experiences with GPT-Live-1 in the API}. 10 September 2026. \url{https://openai.com/index/introducing-gpt-live-1-in-the-api/} \\
S12 & OpenAI. \emph{Introducing the Agents API}. 10 September 2026. \url{https://openai.com/index/introducing-the-agents-api/} \\
S13 & Intern-NCP Team. \emph{NCP-ArchPreview Technical Report}. arXiv:2609.10715, 9 September 2026. \url{https://arxiv.org/abs/2609.10715} \\
S14 & Moshkov, I. et al. \emph{An Open Recipe for IMO Gold}. arXiv:2609.10712, 9 September 2026. \url{https://arxiv.org/abs/2609.10712} \\
S15 & He, J. et al. \emph{ZGCM-1}. arXiv:2609.13356, 11 September 2026. \url{https://arxiv.org/abs/2609.13356} \\
S16 & Zheng, P. et al. \emph{SWE-Bench Pro Verified}. arXiv:2609.08149, 8 September 2026. \url{https://arxiv.org/abs/2609.08149} \\
S17 & Anthropic. \emph{Detecting and countering misuse of AI: September 2026}. 10 September 2026. \url{https://www.anthropic.com/threat-intelligence-report-september-2026} \\
S18 & OpenAI. \emph{Paul Christiano joins OpenAI Foundation Board}. 9 September 2026. \url{https://openai.com/index/paul-christiano-joins-openai-foundation-board/} \\
S19 & Fortune. \emph{Saudi AI giant races to fund data center expansion}. 9 September 2026. \url{https://fortune.com/2026/09/09/saudis-humain-turns-to-outside-investment-as-the-kingdom-reins-in-fiscal-spending/} \\
S20 & NVIDIA. \emph{NVIDIA to Acquire Hugging Face}. 3 September 2026. \url{https://blogs.nvidia.com/blog/nvidia-to-acquire-hugging-face/} \\
\bottomrule
\end{tabularx}}

\vspace{0.8em}
{\footnotesize\color{brandgrey}
Compiled \briefdate. Every figure in this brief was read from the linked primary
source on that date. Corrections: \texttt{\briefauthor}.}

\end{document}
